\documentclass[12pt,a4paper]{article}

% =================================================
% PACOTES
% =================================================
\usepackage[utf8]{inputenc}
\usepackage[T1]{fontenc}
\usepackage[brazil]{babel}
\usepackage{helvet}             % Carrega a fonte Arial / Helvetica
\renewcommand{\familydefault}{\sfdefault} % Define a fonte sans-serif (Arial) para todo o documento
\usepackage{geometry}
\usepackage{array}
\usepackage{tabularx}
\usepackage{longtable}
\usepackage{booktabs}
\usepackage{enumitem}
\usepackage{ragged2e}
\usepackage{setspace}
\usepackage{titlesec}
\usepackage{hyperref}
\usepackage{graphicx}
\usepackage{float}
\usepackage{xurl}

% =================================================
% PÁGINA
% =================================================
\geometry{
    top=3cm,
    bottom=2cm,
    left=3cm,
    right=2cm
}

\onehalfspacing
\setlength{\parindent}{1.25cm}
\setlength{\parskip}{0pt}

\hypersetup{
    colorlinks=true,
    linkcolor=black,
    urlcolor=black
}

% =================================================
% TABELAS
% =================================================
\renewcommand{\arraystretch}{1.3}

\newcolumntype{Y}{>{\RaggedRight\arraybackslash}X}
\newcolumntype{P}[1]{>{\RaggedRight\arraybackslash}p{#1}}

% =================================================
% TÍTULOS
% =================================================
\titleformat{\section}
  {\normalfont\bfseries}
  {\thesection.}
  {0.5em}
  {}

\titleformat{\subsection}
  {\normalfont\bfseries}
  {\thesubsection}
  {0.5em}
  {}

% =================================================
% DOCUMENTO
% =================================================
\begin{document}

% =================================================
% CAPA
% =================================================

\begin{titlepage}
\begin{center}

{\large\textbf{INSTITUTO FEDERAL DE ALAGOAS}}\\
{\large\textbf{CAMPUS MACEIÓ}}\\
{\large\textbf{BACHARELADO EM SISTEMAS DE INFORMAÇÃO}}

\vfill

{\Large\textbf{TERMO DE ABERTURA DE PROJETO}}\\[0.3cm]
{\Large\textbf{(TAP)}}

\vfill

{\large\textbf{SOLIGEST}}\\
{\large Sistema de Gestão Solidária}

\vfill

\textbf{Equipe}\\[0.2cm]
Emerson Gomes Vanderlei da Silva

\vfill

MACEIÓ/AL\\
2026

\end{center}
\end{titlepage}

% =================================================
% TERMO DE ABERTURA
% =================================================

\section*{TERMO DE ABERTURA DO PROJETO (TAP)}

\begin{tabularx}{\textwidth}{P{4.2cm}Y}
\textbf{Nome do Projeto:} &
SoliGest -- Sistema de Gestão Solidária \\

\textbf{Gestor do Projeto:} &
Emerson Gomes \\

\textbf{Prazo do Projeto:} &
12 semanas (4 Sprints) \\

\textbf{Versão do Documento:} &
1.0 \\
\end{tabularx}

\vspace{0.5cm}

\begin{center}
\textbf{Controle de Versões}
\end{center}

\begin{tabularx}{\textwidth}{P{1.5cm}P{5cm}P{3cm}Y}
\toprule
\textbf{Versão} & \textbf{Descrição} & \textbf{Data} & \textbf{Responsável}\\
\midrule
1.0 & Termo de Abertura do Projeto -- SoliGest & A definir & Emerson Gomes\\
\bottomrule
\end{tabularx}

% =================================================
% INTRODUÇÃO
% =================================================

\section*{INTRODUÇÃO}

O SoliGest -- Sistema de Gestão Solidária é uma proposta de aplicação
Web voltada à gestão de instituições do terceiro setor, especialmente
organizações não governamentais (ONGs).

O sistema tem como finalidade centralizar processos relacionados ao
cadastro e controle de voluntários, à triagem de doações financeiras e
de suprimentos, ao agendamento de atendimentos sociais e à geração de
relatórios gerenciais para apoio à tomada de decisão.

Muitas ONGs de pequeno e médio porte ainda dependem de controles
manuais ou de planilhas descentralizadas. Esse cenário pode acarretar
perda de informações, dificuldade no acompanhamento de doações e falta
de transparência com os doadores.

Nesse contexto, o SoliGest busca automatizar esses processos, promovendo
governança de dados, rastreabilidade das informações e ganho
operacional.

\textbf{Principais problemas identificados:}

\begin{itemize}
    \item Dependência de controles manuais;
    \item Utilização de planilhas descentralizadas;
    \item Possibilidade de perda de informações;
    \item Dificuldade no acompanhamento de doações;
    \item Falta de transparência com doadores;
    \item Dificuldade de centralização das informações relacionadas aos
    voluntários e atendimentos sociais.
\end{itemize}

% =================================================
% OBJETIVOS
% =================================================

\section*{OBJETIVOS}

\subsection*{1. Objetivo Geral}

Desenvolver uma aplicação Web capaz de centralizar a gestão de uma
instituição do terceiro setor, permitindo o cadastro e controle de
voluntários, a triagem de doações financeiras e de suprimentos, o
agendamento de atendimentos sociais e a geração de relatórios gerenciais
para tomada de decisão.

\subsection*{2. Objetivos Específicos}

\begin{itemize}
    \item Centralizar o cadastro e controle de voluntários;
    \item Organizar a triagem de doações financeiras;
    \item Organizar a triagem de doações de suprimentos;
    \item Permitir o agendamento de atendimentos sociais;
    \item Disponibilizar relatórios gerenciais;
    \item Promover governança e rastreabilidade dos dados;
    \item Reduzir a dependência de controles manuais e planilhas
    descentralizadas;
    \item Melhorar a organização e o acompanhamento das informações da
    instituição.
\end{itemize}

% =================================================
% PRODUTO E REQUISITOS
% =================================================

\section*{PRODUTO E REQUISITOS}

O produto proposto é um sistema Web denominado SoliGest -- Sistema de
Gestão Solidária, destinado à centralização dos processos de gestão de
uma instituição do terceiro setor.

\textbf{Funcionalidades previstas:}

\begin{enumerate}
    \item Cadastro e controle de voluntários;
    \item Triagem e controle de doações financeiras;
    \item Triagem e controle de doações de suprimentos;
    \item Agendamento de atendimentos sociais;
    \item Geração de relatórios gerenciais;
    \item Centralização das informações da instituição.
\end{enumerate}

\textbf{Requisitos gerais identificados:}

\begin{itemize}
    \item Sistema Web;
    \item Banco de dados relacional;
    \item Interface reativa e dinâmica;
    \item Controle e armazenamento estruturado das informações;
    \item Geração de relatórios gerenciais.
\end{itemize}

% =================================================
% STAKEHOLDERS
% =================================================

\section*{USUÁRIOS / PARTES INTERESSADAS (STAKEHOLDERS)}

O documento original do projeto não apresenta uma relação nominal de
stakeholders nem uma classificação formal de interesse e influência.
Essas informações deverão ser definidas posteriormente.

\begin{center}
\begin{tabularx}{\textwidth}{P{4cm}Y P{2cm}}
\toprule
\textbf{Stakeholder} & \textbf{Interesse} & \textbf{Influência}\\
\midrule
Gestores da ONG & Gestão e acompanhamento das atividades & A definir\\
Voluntários & Cadastro e participação nas atividades da instituição & A definir\\
Doadores & Transparência e acompanhamento das doações & A definir\\
Beneficiários/Atendidos & Agendamento e acompanhamento dos atendimentos & A definir\\
Equipe do projeto & Desenvolvimento e manutenção do sistema & A definir\\
\bottomrule
\end{tabularx}
\end{center}

% =================================================
% PRECIFICAÇÃO
% =================================================

\section*{SIMULAÇÃO DE PRECIFICAÇÃO DO SISTEMA}

O documento original não apresenta uma estimativa de custos de
desenvolvimento ou infraestrutura. Dessa forma, os valores abaixo
permanecem como itens a serem definidos caso a disciplina exija a
simulação de precificação.

\begin{center}
\begin{tabularx}{\textwidth}{P{4cm}P{3cm}P{3cm}Y}
\toprule
\textbf{Função/Item} & \textbf{Quantidade} &
\textbf{Valor Estimado} & \textbf{Observação}\\
\midrule
Desenvolvimento & A definir & A definir & Projeto acadêmico\\
Infraestrutura & A definir & A definir & Ambiente de execução\\
Domínio/Hospedagem & A definir & A definir & Caso necessário\\
\bottomrule
\end{tabularx}
\end{center}

\textbf{Custo Total Estimado:} A definir.

% =================================================
% RISCOS
% =================================================

\section*{RISCOS DO PROJETO E ESTRATÉGIAS DE MITIGAÇÃO}

A governança e a gestão de riscos do projeto consideram a integridade
dos dados, a segurança, o alinhamento arquitetural e a mitigação de
falhas durante a modelagem e operação.

\begin{longtable}{P{0.8cm}P{3.1cm}P{2cm}P{1.8cm}P{6cm}}
\toprule
\textbf{ID} & \textbf{Risco} & \textbf{Probabilidade} &
\textbf{Impacto} & \textbf{Estratégia de Mitigação}\\
\midrule
R1 &
Redundância ou inconsistência nos cadastros de doadores e voluntários &
Média & Alto &
Implementação de regras de validação no Back-End e aplicação de
restrições de integridade no SGBD.\\

R2 &
Vazamento ou acesso indevido a dados sensíveis de
voluntários/atendidos &
Baixa & Crítico &
Aplicação das diretrizes da LGPD, controle de acesso e uso de
criptografia para senhas.\\

R3 &
Modelagem inadequada do banco de dados gerando retrabalho &
Média & Alto &
Validação prévia do Diagrama Entidade-Relacionamento antes do início
da codificação.\\

R4 &
Aumento descontrolado de requisitos &
Média & Médio &
Definição clara do Produto Mínimo Viável (MVP).\\

R5 &
Sobrecarga de papéis e potencial atraso nas entregas devido ao
desenvolvimento solo &
Média & Alto &
Delimitação rigorosa do escopo, focado em um Produto Mínimo Viável
(MVP) funcional.\\
\bottomrule
\end{longtable}

% =================================================
% PADRONIZAÇÃO
% =================================================

\section*{PADRONIZAÇÃO DE NOMENCLATURAS}

Para garantir organização, legibilidade e facilidade de manutenção do
sistema, serão adotadas convenções de nomenclatura padronizadas durante
o desenvolvimento do projeto.

No código do sistema, será adotado um padrão consistente para nomes de
variáveis, funções, classes e componentes. A convenção específica ainda
deverá ser definida de acordo com as práticas adotadas durante a
implementação.

No banco de dados, também deverá ser estabelecido um padrão de
nomenclatura para tabelas e colunas, garantindo consistência e
facilidade de manutenção.

O padrão de mensagens de commit no repositório Git também deverá ser
definido para manter o histórico do projeto organizado.

% =================================================
% PESQUISA DE SISTEMAS EXISTENTES
% =================================================

\section*{PESQUISA SOBRE SISTEMAS DE GESTÃO PARA ONGs}

O documento original do projeto não apresenta uma pesquisa comparativa
com sistemas existentes. Dessa forma, esta seção deverá ser
complementada caso a disciplina exija uma análise de soluções
semelhantes.

O objetivo dessa pesquisa deverá ser identificar funcionalidades
presentes em sistemas de gestão de organizações do terceiro setor que
possam servir como referência para o desenvolvimento do SoliGest.

% =================================================
% REQUISITOS FUNCIONAIS
% =================================================

\section*{REQUISITOS DO SISTEMA}

\subsection*{Requisitos Funcionais}

Os requisitos funcionais descrevem as funcionalidades que o sistema
deverá oferecer aos seus usuários.

\textbf{RF-01 -- Cadastro de voluntários}

O sistema deverá permitir o cadastro e o controle de voluntários da
instituição.

\textbf{RF-02 -- Gestão de doações financeiras}

O sistema deverá permitir o registro e o acompanhamento de doações
financeiras.

\textbf{RF-03 -- Gestão de doações de suprimentos}

O sistema deverá permitir o registro e a triagem de doações de
suprimentos.

\textbf{RF-04 -- Agendamento de atendimentos sociais}

O sistema deverá permitir o agendamento de atendimentos sociais.

\textbf{RF-05 -- Relatórios gerenciais}

O sistema deverá permitir a geração de relatórios gerenciais para
auxiliar a tomada de decisão.

\textbf{RF-06 -- Centralização das informações}

O sistema deverá centralizar as informações relacionadas aos
voluntários, doações e atendimentos sociais.

% =================================================
% REQUISITOS NÃO FUNCIONAIS
% =================================================

\subsection*{Requisitos Não Funcionais}

Os requisitos não funcionais descrevem características de qualidade e
restrições técnicas do sistema.

\textbf{RNF-01 -- Plataforma Web}

O sistema deverá ser disponibilizado por meio de uma aplicação Web.

\textbf{RNF-02 -- Usabilidade}

A aplicação deverá apresentar interfaces reativas, componentizadas e
dinâmicas.

\textbf{RNF-03 -- Banco de Dados}

O sistema deverá utilizar um SGBD relacional para armazenamento
estruturado das informações.

\textbf{RNF-04 -- Segurança}

O sistema deverá considerar mecanismos de segurança e proteção dos
dados, incluindo as diretrizes aplicáveis da LGPD.

\textbf{RNF-05 -- Testabilidade}

O projeto deverá contemplar testes unitários e testes de ponta a ponta.

% =================================================
% TECNOLOGIAS
% =================================================

\section*{TECNOLOGIAS PROPOSTAS PARA O PROJETO}

As tecnologias selecionadas para o desenvolvimento do SoliGest são
apresentadas a seguir.

\subsection*{Back-End}

\begin{itemize}
    \item \textbf{Node.js} -- utilizado no desenvolvimento do Back-End,
    devido ao desempenho, ecossistema robusto e facilidade na construção
    de APIs RESTful.
\end{itemize}

\subsection*{Front-End}

\begin{itemize}
    \item \textbf{React.js} -- utilizado para criação de interfaces
    reativas, componentizadas e dinâmicas.
\end{itemize}

\subsection*{Banco de Dados}

\begin{itemize}
    \item \textbf{MySQL} -- SGBD consolidado, de código aberto, com
    amplo suporte e integração com Node.js.
\end{itemize}

\subsection*{Prototipação e Design}

\begin{itemize}
    \item \textbf{Figma} -- utilizado para criação de wireframes de
    alta fidelidade e design system responsivo.
\end{itemize}

\subsection*{Gerenciamento de Projeto}

\begin{itemize}
    \item \textbf{Trello} -- utilizado para organização do fluxo de
    trabalho utilizando metodologia ágil.
\end{itemize}

\subsection*{Controle de Versão}

\begin{itemize}
    \item \textbf{GitHub} -- utilizado para hospedagem do repositório
    remoto, versionamento Git, controle de commits e rastreabilidade
    do código.
\end{itemize}

\subsection*{Ambiente de Desenvolvimento}

\begin{itemize}
    \item \textbf{VS Code} -- ambiente de desenvolvimento leve,
    gratuito, extensível e integrado ao Git.
\end{itemize}

\subsection*{Testes}

\begin{itemize}
    \item \textbf{Postman} -- validação de endpoints, automação de
    coleções e testes de integração REST;
    \item \textbf{Jest} -- testes unitários;
    \item \textbf{Cypress} -- testes de ponta a ponta (E2E).
\end{itemize}

% =================================================
% ARQUITETURA
% =================================================

\section*{ARQUITETURA DO SISTEMA}

\subsection*{1. Visão Geral da Arquitetura}

O documento original define as principais tecnologias do sistema, mas
não apresenta uma arquitetura detalhada em camadas. A arquitetura
definitiva deverá ser especificada durante as etapas de modelagem e
desenvolvimento.

A solução será desenvolvida como uma aplicação Web, utilizando React.js
no Front-End, Node.js no Back-End e MySQL como sistema gerenciador de
banco de dados.

\subsection*{2. Camada de Apresentação (Front-End)}

A camada de apresentação será desenvolvida utilizando React.js,
responsável pela criação das interfaces reativas, componentizadas e
dinâmicas do sistema.

\subsection*{3. Camada de Aplicação (Back-End)}

A camada de aplicação utilizará Node.js para implementar a lógica de
negócio e as APIs RESTful responsáveis pela comunicação entre o
Front-End e o banco de dados.

\subsection*{4. Camada de Dados}

A camada de dados será implementada utilizando MySQL, responsável pelo
armazenamento estruturado das informações relacionadas aos voluntários,
doações e atendimentos.

\subsection*{5. Segurança}

A segurança deverá considerar controle de acesso e proteção dos dados,
especialmente das informações relacionadas aos voluntários e
beneficiários/atendidos, observando as diretrizes da LGPD.

% =================================================
% CRONOGRAMA / MARCOS
% =================================================

\section*{MARCOS DO PROJETO}

O cronograma foi estruturado em quatro Sprints principais, totalizando
12 semanas.

\begin{itemize}
    \item \textbf{Sprint 1 -- Semana 1--3:} Levantamento, requisitos,
    modelagem e prototipação;
    \item \textbf{Sprint 2 -- Semana 4--6:} Modelagem do SGBD,
    configuração do ambiente e desenvolvimento da API base (CRUDs);
    \item \textbf{Sprint 3 -- Semana 7--9:} Desenvolvimento do
    Front-End, integração com Back-End e implementação das regras
    negociais;
    \item \textbf{Sprint 4 -- Semana 10--12:} Execução de testes (QA),
    ajustes, documentação final e apresentação.
\end{itemize}

% =================================================
% EQUIPE
% =================================================

\section*{EQUIPE}

\begin{itemize}
    \item \textbf{Emerson Gomes} -- responsável pelo desenvolvimento
    do projeto.
\end{itemize}

Como o projeto é desenvolvido individualmente, as demais funções
necessárias ao projeto serão desempenhadas pelo próprio integrante,
conforme a necessidade de cada etapa.

\end{document}
